% Figure: the power-mean ladder. For two fixed positive numbers a and b
% the power mean M_p is monotone non-decreasing in p, passing through the
% harmonic, geometric, arithmetic, and quadratic means.
\begin{figure}[htbp]
\centering
\begin{tikzpicture}
\begin{axis}[
  width=11cm, height=6.5cm,
  xlabel={exponent $p$},
  ylabel={power mean $M_{p}(2,8)$},
  xmin=-4, xmax=4,
  ymin=2.5, ymax=8.5,
  grid=both,
  grid style={engblack!12},
  tick label style={font=\footnotesize},
  label style={font=\small},
  legend style={font=\footnotesize, at={(0.98,0.04)}, anchor=south east},
]
% M_p(2,8) = ((2^p + 8^p)/2)^(1/p). Sampled values (p, M_p).
\addplot[engblue, very thick, smooth] coordinates {
  (-4, 2.262) (-3, 2.397) (-2, 2.659) (-1.5, 2.857) (-1, 3.200)
  (-0.5, 3.660) (-0.25, 3.823) (0, 4.000) (0.25, 4.166) (0.5, 4.317)
  (1, 5.000) (1.5, 5.392) (2, 5.831) (3, 6.604) (4, 7.110)
};
% Marked classical means.
\addplot[only marks, mark=*, engred, mark size=2.2pt] coordinates {
  (-1, 3.200) (0, 4.000) (1, 5.000) (2, 5.831)
};
\node[font=\footnotesize, anchor=north] at (axis cs:-1,3.20) {HM};
\node[font=\footnotesize, anchor=north] at (axis cs:0,4.00) {GM};
\node[font=\footnotesize, anchor=south] at (axis cs:1,5.00) {AM};
\node[font=\footnotesize, anchor=south] at (axis cs:2,5.83) {QM};
% asymptotes min and max
\addplot[engblack!50, dashed] coordinates {(-4,2)(4,2)};
\addplot[engblack!50, dashed] coordinates {(-4,8)(4,8)};
\node[font=\footnotesize, anchor=west] at (axis cs:-3.8,2.25) {$\min=2$};
\node[font=\footnotesize, anchor=west] at (axis cs:-3.8,7.75) {$\max=8$};
\legend{{$M_{p}(2,8)$}, {classical means}}
\end{axis}
\end{tikzpicture}
\caption[The power-mean ladder]{The power mean
$M_{p}(a,b)=\bigl((a^{p}+b^{p})/2\bigr)^{1/p}$ is non-decreasing in
$p$ for fixed positive $a,b$, shown here for $a=2$, $b=8$. The harmonic
mean ($p=-1$), geometric mean ($p\to 0$), arithmetic mean ($p=1$), and
quadratic (root-mean-square) mean ($p=2$) sit on the same curve in
increasing order, which is the chain $\mathrm{HM}\le\mathrm{GM}\le
\mathrm{AM}\le\mathrm{QM}$. The curve tends to $\min(a,b)$ as
$p\to-\infty$ and to $\max(a,b)$ as $p\to+\infty$. AM--GM is the single
step from $p=0$ to $p=1$.}
\label{fig:vol02:ch01:power-mean}
\end{figure}
